\section{AI交互平台与Geant4到TCAD的传递}
\label{cr:platform}
\subsection{自然语言、结构化设计与执行边界}
AI在项目中辅助阅读文献、组织假设、生成和检查脚本、定位失败、提取结果及
撰写解释。可以被复验的交互接口则是限定中文模板：请求明确列出栅长、
粒子、能量、方向、重复数和范围，解析后得到结构化工况清单。
每个工况包含几何条件、粒子条件、种子、配对位值及训练或测试身份。
接口拒绝模糊尺寸、未支持角度和超出冻结范围的要求，避免把自然语言中的
省略或歧义直接传给求解器。

v1已有中文请求生成完整矩阵并进入真实执行链路的记录。
当前R50演示入口提供\texttt{preview}、\texttt{audit}和\texttt{locate}：
分别预览映射、检查接受与拒绝行为、定位指定事件的证据路径。
它只说明请求与已冻结R50设计如何对应，不自动提交新求解。
$50\,\Omega$来自版本化物理修订，而非解析器从句子中猜测出的参数。
自然语言成功解析是软件层的检查，实际模型正确性仍由后续物理与数值证据判定。

\begin{figure}[htbp]
\centering
% Conceptual workflow, not a measured result or screenshot.
\resizebox{\linewidth}{!}{%
\begin{tikzpicture}[
 block/.style={draw=signalblue,rounded corners=2pt,fill=softgray,
 text width=4.2cm,minimum height=1.5cm,align=center,font=\small},
 arrow/.style={-{Latex[length=2.4mm]},thick,draw=signalblue},
 node distance=0.7cm and 0.6cm]
\node[block] (request) {限定中文请求\\尺寸、粒子、角度\\独立事件数};
\node[block,right=of request] (matrix) {受约束的工况清单\\单位、范围、版本与种子};
\node[block,right=of matrix] (transport) {Geant4 硅内电离沉积\\$E_{\rm ion}\ \rightarrow\ Q_{\rm dep}$\\有限径迹及单位检查};
\node[block,below=of transport] (tcad) {TCAD 器件与电路求解\\等效径迹与受控初态\\读出瞬态与端口检查};
\node[block,left=of tcad] (metric) {匹配中性参考与同位对照\\$D,\ M,\ \Delta Q,\ \Delta M$};
\node[block,left=of metric] (model) {按工况分组的响应拟合\\冻结预测后检验未见响应};
\draw[arrow] (request)--(matrix);
\draw[arrow] (matrix)--(transport);
\draw[arrow] (transport)--(tcad);
\draw[arrow] (tcad)--(metric);
\draw[arrow] (metric)--(model);
\node[below=0.40cm of metric,align=center,font=\small,text width=14.5cm]
 {每一层绑定输入版本、运行记录和来源哈希；科学质量门槛与执行状态分别记录。};
\end{tikzpicture}}

\caption{AI辅助联合评估的工作流程示意。自然语言请求经结构化审查与冻结后，
连接粒子输运、等效源传递、器件求解和独立评价；此图为流程示意，不代表未完成阶段已有结果。}
\label{cr:workflow}
\end{figure}

\subsection{输运记录与等效源转换}
Geant4在硅与真空组成的简化几何中运行电磁option4及步长限制，
记录每一步的位置、时间、总沉积能、非电离沉积能和轨迹身份。
主矩阵采用$5\,\mathrm{MeV}$的$\alpha$粒子与质子，方向为0°和30°，
表面入射位置与各栅长模型保持约定关系。30°径迹可能从侧面退出，
因此不能机械地把硅深度除以角度余弦视为所有工况的有效径迹长度。
实际几何交点及源文件共同决定传递参数。

只把硅内电离能量计入电子空穴对产生。设逐步电离沉积总量为
$E_{\mathrm{ion}}$，使用硅中平均$3.6\,\mathrm{eV}$产生一对载流子的近似，
得到
\begin{equation}
N_{\mathrm{eh}}=\frac{10^6E_{\mathrm{ion}}[\mathrm{MeV}]}{3.6},\qquad
Q_{\mathrm{dep}}=qN_{\mathrm{eh}},\qquad
\mathrm{LET}_{\mathrm{eq}}=\frac{Q_{\mathrm{dep}}}{\ell}.
\label{cr:transfer}
\end{equation}
式中$\ell$是约定的硅内等效直径迹长度。TCAD以半径$0.03\um$的
Gaussian HeavyIon源注入等效电荷，并使用一致的二维有效厚度和电荷单位。
这一方法保留事件的电离总量，便于小规模参数扫描；逐步输运CSV仍单独保留，
以便识别总量近似忽略的纵向、横向和时间分布信息。

已有TCAD与Geant4的SRAM评估研究提供了联合方法的先例\cite{chen2011}；
硅二极管中重离子、聚焦X射线和脉冲激光的瞬态比较，则提示不同电荷产生方式
需要分别处理\cite{ryder2021}。这些文献不能直接校准本项目的DRAM模型。
本项目采用总量等效直径迹而非完整三维产生率网格，
故需要分别检查转换算术与TCAD实际源积分。前者核对单位和总电荷，
后者检查离散网格、径迹截断和时间积分之后是否仍近似注入目标电荷。
二者通过并不意味着空间分布误差消失。

\subsection{可追溯执行与数据流}
执行流程将设计、输入deck、Geant4事件三文件、求解日志及输出归档绑定到唯一
run ID，并用SHA-256检查内容一致性。续算只允许使用同几何、同电路、同网格和
已核对合同的保存点；改变网格或参考结构后重新建立相应初态。
主运行器负责归档，分析器从固定原始输出提取曲线，图表再引用这些曲线和摘要，
形成“请求—输入—源—求解—指标—图”的连续来源关系。

这种设计尤其适合AI辅助操作。AI能够快速生成大量命令，也可能误读旧摘要、
把同名文件当作同一内容，或把后台退出当成远端求解停止。
因此流程把可机器检查的版本与完成条件写入入口，而不是依赖对话记忆。
报告中的数据图属于原始数值后处理；真实Sentaurus窗口截图和原生场图另行标注，
防止用美观的曲线图替代软件求解证据。
